% hessian.tex -- the RHF and CASSCF analytic nuclear Hessians: their design
% derivations, the parts the module READMEs, sources, tests and reference
% generators under src/, test/ and tools/ cite.
%
% These are the design notes of 31 October 2026, written before src/hessian
% and the modules beneath it, and they keep that note's voice: "the tree
% lacks" means it lacked it then. Each section ends with what has landed
% since, and the module READMEs are the record of the code as it stands;
% where the two disagree, the code and its READMEs win. Equation labels are
% the notes' own, so a README's "eq. (sa)" is \eqref{eq:sa} here, and a
% section, algorithm or equation is cited by its label without the prefix
% ("section roots", "Algorithm ss"); the RHF part's labels that collide with
% the CASSCF part's carry an rhf- prefix (sec:rhf-notation, eq:rhf-energy,
% sec:rhf-gradient, eq:rhf-gradient, sec:rhf-df, sec:rhf-algorithm,
% sec:rhf-checks, sec:rhf-plan). Code names are the tree's current ones.
% The CASSCF notes' inventory of what then existed is left out: the module
% READMEs are that record now.
%
% Build: pdflatex hessian.tex, twice (cross-references).
\documentclass[11pt]{article}

\usepackage[T1]{fontenc}
\usepackage{lmodern}
\usepackage{microtype}
\usepackage[margin=2.5cm]{geometry}
\usepackage{amsmath,amssymb}
\usepackage{booktabs}
\usepackage{algorithm}
\usepackage[noend]{algpseudocode}
\usepackage[hidelinks]{hyperref}

\newcommand{\code}[1]{\texttt{#1}}
\newcommand{\eps}{\varepsilon}
\newcommand{\tr}{\operatorname{tr}}
\newcommand{\asym}[2]{(#1\Vert#2)}
\newcommand{\Fsk}[1]{F^{(#1)}}
\newcommand{\G}[1]{\Gamma^{(#1)}}
\newcommand{\Lag}{\mathcal L}
\newcommand{\Exy}{E^{xy}}
\newcommand{\sym}{\operatorname{sym}}
\newcommand{\Dsk}[1]{D_S^{#1}}
\algrenewcommand\algorithmicrequire{\textbf{Input:}}
\algrenewcommand\algorithmicensure{\textbf{Output:}}

\title{The RHF and CASSCF analytic nuclear Hessians\\
\large design notes: derivation, coupled-perturbed equations, density
fitting and the state average}
\author{free-breakfast design notes}
\date{31 October 2026; brought into the tree 8 November 2026}

\begin{document}
\sloppy
\maketitle

\begin{abstract}
Part~I is the closed-shell RHF Hessian: the derivation in the vocabulary of
the tree's RHF gradient, two equivalent forms of the result (one manifestly
symmetric, one made only of AO matrices and Fock builds), the CPHF
equations, the skeleton Fock derivatives every right-hand side needs, the
density-fitted variant, the algorithm, its checks and a build order. By
Wigner's $2n+1$ rule no second-order response is needed. Part~II is the
CASSCF Hessian, state-specific and state-averaged: the second-order
connection, the CP-CASSCF right-hand side and its $3N$ solves with the
$Z$-vector's operator, the state average's Lagrangian, whose second
derivative needs the gradient's own $Z$ and a third derivative of the
averaged energy contracted once with it, the rotations among the averaged
roots, the explicit part over density fitting, the algorithm, the
conditions it needs, its checks and a build order.
\end{abstract}

\tableofcontents

%======================================================================
\part{RHF}

\section{Notation and the energy}
\label{sec:rhf-notation}

Orbitals are real and the reference is closed-shell RHF with $n_o$ doubly
occupied orbitals. AO indices are $\mu\nu\lambda\sigma$, occupied MOs
$ijkl$, virtual MOs $abcd$, and general MOs $pqrs$. $C$ is the MO
coefficient matrix, $C_o$ and $C_v$ its occupied and virtual columns, and
$\eps_p$ the canonical orbital energies. Nuclear coordinates are $x$ and
$y$; each is one Cartesian component of one nucleus, $x = (A, k)$, and there
are $3N$ of them for $N$ atoms.

The density and energy-weighted density are
\begin{equation}
  D = 2\,C_o C_o^{\mathsf T}, \qquad
  W = 2\,C_o\,\eps_o\,C_o^{\mathsf T} = \tfrac12\,D F D,
  \label{eq:dw}
\end{equation}
the second form of $W$ holding for any $S$-orthonormal $C$ that
diagonalizes the occupied block of $F$ (since $C_o^{\mathsf T} F C_o =
\eps_o$). It is the form that differentiates cleanly.

With the antisymmetrized integral
\begin{equation}
  \asym{\mu\nu}{\lambda\sigma} = (\mu\nu|\lambda\sigma)
  - \tfrac12(\mu\lambda|\nu\sigma),
\end{equation}
which is symmetric under $(\mu\nu)\leftrightarrow(\lambda\sigma)$ for real
functions, the two-electron operator on a symmetric density $P$ is
\begin{equation}
  G(P)_{\mu\nu} = \sum_{\lambda\sigma} P_{\lambda\sigma}
  \asym{\mu\nu}{\lambda\sigma} = J(P)_{\mu\nu} - \tfrac12 K(P)_{\mu\nu},
\end{equation}
and $G$ is a symmetric bilinear form: $\tr[P\,G(Q)] = \tr[Q\,G(P)]$. The
Fock matrix is $F = h + G(D)$ and the energy is
\begin{equation}
  E = \sum_{\mu\nu} D_{\mu\nu} h_{\mu\nu}
  + \tfrac12 \sum_{\mu\nu\lambda\sigma} D_{\mu\nu} D_{\lambda\sigma}
  \asym{\mu\nu}{\lambda\sigma} + V_{nn}.
  \label{eq:rhf-energy}
\end{equation}
The SCF conditions are $F_{ai} = 0$ and $C^{\mathsf T} S C = 1$, at every
geometry.

\paragraph{Derivative integrals.} A superscript on an integral means its
derivative at fixed AO coefficients, the basis functions moving with their
nuclei: $h^x_{\mu\nu} = \partial h_{\mu\nu}/\partial x$, $S^x$, and
$(\mu\nu|\lambda\sigma)^x$; likewise $h^{xy}$, $S^{xy}$ and
$(\mu\nu|\lambda\sigma)^{xy}$. $h^x$ includes the derivative of the
nuclear-attraction \emph{operator} as well as of the basis functions. MO
transforms with the unperturbed coefficients are written $S^x_{pq} =
C_p^{\mathsf T} S^x C_q$, and similarly for any AO matrix.

\paragraph{Skeleton Fock derivative.} The derivative of $F$ at fixed $D$,
\begin{equation}
  \Fsk{x}_{\mu\nu} = h^x_{\mu\nu}
  + \sum_{\lambda\sigma} D_{\lambda\sigma}
  \asym{\mu\nu}{\lambda\sigma}^x,
  \label{eq:skeleton}
\end{equation}
is the central intermediate. There are $3N$ of them, each an AO matrix.
The Hessian is an explicit part, the second-derivative integrals contracted
with the converged density, plus a response part that needs the
first-order orbital response, one coupled-perturbed Hartree--Fock (CPHF)
solve per nuclear coordinate; $\Fsk{x}$ enters both, and every CPHF
right-hand side.

%======================================================================
\section{The gradient, as the tree computes it}
\label{sec:rhf-gradient}

Differentiating \eqref{eq:rhf-energy} and using the SCF conditions gives the
familiar
\begin{equation}
  E^x = \sum D_{\mu\nu} h^x_{\mu\nu}
  + \tfrac12 \sum D_{\mu\nu} D_{\lambda\sigma}
    \asym{\mu\nu}{\lambda\sigma}^x
  - \sum W_{\mu\nu} S^x_{\mu\nu} + V_{nn}^x,
  \label{eq:rhf-gradient}
\end{equation}
which is exactly \code{gradient::rhf}'s term list
(\code{src/gradient/README.md}): \code{kinetic\_gradient} and
\code{potential\_gradient} at $D$, \code{overlap\_gradient} at $-W$,
\code{coulomb\_gradient} on $D$ and $-\frac12$\,\code{exchange\_gradient}
on $D$, and the nuclear repulsion. The orbital response is absent because
$E$ is stationary in the occupied--virtual rotations, and the normalization
response is carried by $-W S^x$.

The Hessian is the derivative of \eqref{eq:rhf-gradient}. Equation
\eqref{eq:rhf-gradient} holds at \emph{every} geometry provided $D$ and $W$
are the converged ones there, so differentiating it is legitimate, and two
kinds of term result: those where $y$ hits an integral (the \emph{explicit}
part) and those where $y$ hits $D$ or $W$ (the \emph{response} part):
\begin{equation}
  \Exy = \Exy_{\mathrm{expl}}
  + \sum D^y_{\mu\nu}\,\Fsk{x}_{\mu\nu}
  - \sum W^y_{\mu\nu}\,S^x_{\mu\nu},
  \label{eq:hess-ao}
\end{equation}
\begin{equation}
  \Exy_{\mathrm{expl}} = \sum D_{\mu\nu} h^{xy}_{\mu\nu}
  + \tfrac12 \sum D_{\mu\nu} D_{\lambda\sigma}
    \asym{\mu\nu}{\lambda\sigma}^{xy}
  - \sum W_{\mu\nu} S^{xy}_{\mu\nu} + V_{nn}^{xy}.
  \label{eq:explicit}
\end{equation}
The middle term of \eqref{eq:hess-ao} collects $D^y h^x$ and the two halves
of $\frac12 (D D)^y \asym{}{}^x$, which are equal by the pair symmetry of
$\asym{}{}$. Here $D^y$ and $W^y$ are \emph{total} derivatives: they need
the orbitals' response.

%======================================================================
\section{The orbital response}
\label{sec:response}

\subsection{Parametrization and the orthonormality constraint}

Expand the perturbed coefficients in the unperturbed ones,
\begin{equation}
  \frac{\partial C_{\mu p}}{\partial y} = \sum_q C_{\mu q}\,U^y_{qp}.
\end{equation}
Differentiating $C^{\mathsf T} S C = 1$ gives
\begin{equation}
  U^y_{pq} + U^y_{qp} + S^y_{pq} = 0.
  \label{eq:ortho}
\end{equation}
The occupied--occupied block of $U^y$ is a rotation among occupied orbitals
plus a symmetric part fixed by \eqref{eq:ortho}. $D$ and $E$ are invariant
to the rotation, so we choose
\begin{equation}
  U^y_{ij} = -\tfrac12 S^y_{ij}
  \label{eq:uoo}
\end{equation}
(the orbitals do \emph{not} stay canonical along $y$; this is why $W$
is better differentiated as $\frac12 DFD$ than as $2C_o\eps_oC_o^{\mathsf T}$).
The virtual--occupied block $U^y_{ai}$ is the unknown, and the
occupied--virtual block follows from \eqref{eq:ortho}:
$U^y_{ia} = -U^y_{ai} - S^y_{ai}$. The virtual--virtual block never
enters a closed-shell quantity.

\subsection{The density derivative}

With \eqref{eq:uoo},
\begin{equation}
  D^y = D_U^y - \Dsk{y}, \qquad
  D_U^y = 2\sum_{ai} U^y_{ai}\,\bigl(C_a C_i^{\mathsf T} + C_i C_a^{\mathsf T}\bigr),
  \qquad
  \Dsk{y} = 2\,C_o\,S^y_{oo}\,C_o^{\mathsf T},
  \label{eq:dy}
\end{equation}
where $S^y_{oo}$ is the $n_o\times n_o$ block of $S^y_{pq}$. $\Dsk{y}$ is
known from the overlap derivative alone; $D_U^y$ needs the CPHF solution.
In code $D_U^y = 2(X + X^{\mathsf T})$ with $X = C_v U^y C_o^{\mathsf T}$:
two GEMMs.

\subsection{The coupled-perturbed Hartree--Fock equations}

Differentiate the condition $F_{ai} = C_a^{\mathsf T} F C_i = 0$. Three
pieces appear: the coefficient derivatives on either side, the skeleton
derivative \eqref{eq:skeleton}, and the Fock change from $D^y$:
\begin{equation}
  0 = \frac{dF_{ai}}{dy}
  = U^y_{ia}\,\eps_i + U^y_{ai}\,\eps_a + \Fsk{y}_{ai} + G(D^y)_{ai}.
\end{equation}
Substituting $U^y_{ia} = -U^y_{ai} - S^y_{ai}$ and \eqref{eq:dy}:
\begin{equation}
  \boxed{\;
  \sum_{bk} M_{ai,bk}\,U^y_{bk} = B^y_{ai}, \qquad
  B^y_{ai} = -\Fsk{y}_{ai} + \eps_i\,S^y_{ai} + G(\Dsk{y})_{ai}
  \;}
  \label{eq:cphf}
\end{equation}
with the electronic Hessian
\begin{equation}
  M_{ai,bk} = (\eps_a - \eps_i)\,\delta_{ab}\delta_{ik}
  + 4(ai|bk) - (ab|ik) - (ak|ib).
  \label{eq:m}
\end{equation}
$M$ is symmetric, and positive definite exactly when the RHF solution is a
minimum under real singlet rotations (it is the real RHF$\to$RHF stability
matrix up to a factor). A CPHF solve that does not converge is therefore a
diagnostic of an unstable reference, not of the solver.

\paragraph{The product never forms $M$.} For a trial $U$,
\begin{equation}
  (MU)_{ai} = (\eps_a - \eps_i)\,U_{ai} + C_a^{\mathsf T}\,G(D_U)\,C_i,
  \label{eq:matvec}
\end{equation}
with $D_U$ built from $U$ as in \eqref{eq:dy}. So one product is two GEMMs
to $D_U$, one J/K build on a symmetric density, and two GEMMs back to the
$ai$ block: the same work as one SCF iteration. The $3N$ right-hand
sides are independent, so the J/K builds batch across them, and the
diagonal $(\eps_a - \eps_i)$ is the obvious preconditioner.

\paragraph{How many solves.} One per coordinate, $3N$, or by translational
invariance $3(N-1)$: the response to moving every nucleus together is a
pure translation of the orbitals, so one atom's three responses are minus
the sum of the others'. Rotational invariance gives three more only at a
stationary point and only with care; it is not worth it. The $2n+1$ rule
says first-order response suffices for a second derivative: no $U^{xy}$ is
ever needed. The $Z$-vector trick does not help here: it removes response
solves when there is one functional to differentiate, and a Hessian has
$3N$ gradient components.

\paragraph{Since written.} The CPHF is \code{response::cphf}
(\code{src/response/cphf}), every right-hand side in one block.

%======================================================================
\section{The Hessian}
\label{sec:hessian}

\subsection{The AO form}

Equation \eqref{eq:hess-ao} with $D^y$ from \eqref{eq:dy} and $W^y$ from
\eqref{eq:dw},
\begin{equation}
  W^y = \tfrac12\bigl(D^y F D + D F^y D + D F D^y\bigr),
  \qquad F^y = \Fsk{y} + G(D^y),
  \label{eq:wy}
\end{equation}
is already a complete algorithm: per coordinate $y$, form $D^y$, one J/K
build for $G(D^y)$, a handful of GEMMs for $W^y$, then $3N$ dot products
$\tr(D^y \Fsk{x}) - \tr(W^y S^x)$ for every $x$. It is not manifestly
symmetric in $x$ and $y$, but it is the exact derivative of an exact
gradient, so its asymmetry measures the CPHF residual. Symmetrizing
$\frac12(\Exy + E^{yx})$ at the end is standard, and the discarded
antisymmetric part is a free convergence diagnostic.

\subsection{The symmetric form}

Expanding \eqref{eq:hess-ao} in MOs, using \eqref{eq:cphf} to eliminate the
$\Fsk{x}_{ai}$ and the bilinear symmetry of $G$ (two terms
$\pm 4\sum_{ai} U^y_{ai}\,G(\Dsk{x})_{ai}$ cancel, one from each of $D^y$
and $W^y$), gives
\begin{equation}
  \boxed{\;
  \begin{aligned}
  \Exy = \Exy_{\mathrm{expl}}
  &- 4\sum_{ai} B^x_{ai}\,U^y_{ai}
   - 2\sum_{ij}\bigl(S^x_{ij}\,\Fsk{y}_{ij} + S^y_{ij}\,\Fsk{x}_{ij}\bigr)\\
  &+ 2\sum_{ij} S^x_{ij}\,S^y_{ij}\,(\eps_i + \eps_j)
   + \tr\bigl[\Dsk{x}\,G(\Dsk{y})\bigr].
  \end{aligned}\;}
  \label{eq:hess-sym}
\end{equation}
Every term is symmetric in $x\leftrightarrow y$; the first because
$\sum_{ai} B^x_{ai} U^y_{ai} = B^{x\mathsf T} M^{-1} B^y$. That term is the
orbital relaxation, and it is negative semi-definite since $M$ is positive
definite: letting the orbitals relax can only lower the curvature. It is
the term that a frozen-orbital Hessian omits and that dominates the error
of one.

The two forms agree to the CPHF tolerance (section~\ref{sec:rhf-checks}).
\eqref{eq:hess-sym} costs fewer J/K builds (the $G(\Dsk{y})$ are needed
for $B^y$ anyway) and needs no $W^y$; \eqref{eq:hess-ao} is shorter to
write over AO matrices. The recommendation is \eqref{eq:hess-sym}, with
\eqref{eq:hess-ao} kept as a test.

\subsection{Nuclear repulsion}

With $\mathbf R = \mathbf R_A - \mathbf R_B$ and $R = |\mathbf R|$, the
off-diagonal blocks are
\begin{equation}
  \frac{\partial^2 V_{nn}}{\partial A_k\,\partial B_l}
  = -Z_A Z_B\,\frac{3R_kR_l - R^2\delta_{kl}}{R^5}
  \qquad (A \neq B),
\end{equation}
and each diagonal block is minus the sum of its row's off-diagonal blocks.

\paragraph{Since written.} The response part, in the symmetric form
\eqref{eq:hess-sym}, is \code{hessian::df\_rhf\_response}
(\code{src/hessian}), and its tests build the AO form \eqref{eq:hess-ao}
from the same $U^x$ as the check this section recommends.

%======================================================================
\section{The explicit part and its integrals}
\label{sec:integrals}

$\Exy_{\mathrm{expl}}$ \eqref{eq:explicit} needs four second-derivative
integral classes, each only ever contracted, never stored:

\begin{center}
\begin{tabular}{@{}lll@{}}
\toprule
class & contracted with & centres differentiated \\
\midrule
$S^{xy}$ & $-W$ & 2 (bra, ket) \\
$T^{xy}$ & $D$ & 2 \\
$V^{xy}$ & $D$ & 2, plus every nucleus $C$ (operator) \\
$(\mu\nu|\lambda\sigma)^{xy}$ & $\frac12 D\otimes D$ (antisymmetrized) & 4 \\
\bottomrule
\end{tabular}
\end{center}

\paragraph{Shape of the output.} A shell pair touches two atoms, so an
overlap or kinetic second derivative fills at most a $6\times6$ block of
the $3N\times3N$ Hessian, and a nuclear-attraction one, whose operator
nucleus $C$ is a third atom, a $9\times9$ block; a quartet touches four, a
$12\times12$ block with 78 unique entries. Translational invariance,
$\sum_A \partial/\partial A_k = 0$ applied to each integral, removes one
centre: a quartet's derivatives need only three centres' worth of distinct
integrals (45 unique second-derivative entries), and an overlap or kinetic
pair's only one. The tree's first-derivative kernels do not use
translational invariance
(\code{src/integrals/one\_electron/gradient/README.md}); for second
derivatives the saving is large enough that they should.

\paragraph{The nuclear-attraction operator.} $V^{xy}$ has three kinds of
term: both derivatives on basis centres, one on a basis centre and one on
the operator's nucleus, and both on the operator's nucleus. The last is
the second derivative of the electrostatic potential at nucleus $C$, the
same integral family as electric-field-gradient integrals. Translational
invariance turns the operator derivatives into basis-centre ones,
$\partial_C = -(\partial_A + \partial_B)$ for the $C$-term of a pair on
$A,B$, which is how most codes avoid a separate kernel.

\paragraph{Angular momentum.} In McMurchie--Davidson form a derivative on
a Gaussian is a sum of functions of $l\pm1$, so a second derivative needs
Hermite expansions to $l+2$ on that centre, and a quartet's to $L+2$ in
total over the two centres differentiated. At f this is the same register
and stack pressure that bounded the f-shell first-derivative and exchange
kernels; the generated kernels' compile time and spill behaviour should be
measured with \code{devtools/kernel-compile-watchdog.sh} before anything
is assumed.

\paragraph{Since written.} The one-electron contractions are
\code{overlap\_hessian}, \code{kinetic\_hessian} and
\code{potential\_hessian} (\code{src/integrals/one\_electron/hessian}),
their diagonal blocks from translational invariance; the fitted
two-electron ones are \code{metric\_hessian} and
\code{three\_center\_hessian}
(\code{src/integrals/eri/density\_fitting/hessian}).

%======================================================================
\section{The skeleton Fock derivatives}
\label{sec:skeleton}

Both forms of the Hessian, and every CPHF right-hand side, need the $3N$
AO matrices $\Fsk{x}$ \eqref{eq:skeleton}. The two-electron part is a
\emph{derivative J/K build}: $J^{(x)}(D)_{\mu\nu} = \sum D_{\lambda\sigma}
(\mu\nu|\lambda\sigma)^x$ and the matching $K^{(x)}$, for every $x$ at
once. This is the kernel the tree most clearly lacks. Its first-derivative
integrals exist (\code{coulomb\_gradient} and \code{exchange\_gradient}
generate $(\mu\nu|\lambda\sigma)^x$), but they contract both index pairs
with $D$ and return $3N$ scalars. A Hessian needs the same integrals
contracted with one $D$, returning $3N$ matrices: $3N\,n_{\mathrm{AO}}^2$
doubles of output, and scatter writes into per-atom slabs.

\paragraph{The one-electron part.} $h^x$ as matrices rather than scalars is
the same change to \code{kinetic\_gradient} and \code{potential\_gradient},
and $S^x$ as matrices to \code{overlap\_gradient}; all three are cheap.

A shell quartet differentiated on the centre of shell $\mu$ contributes to
$J^{(x)}$ through every $(\mu\nu)$ and $(\lambda\sigma)$ element it
touches, so the natural layout is one $n_{\mathrm{AO}}^2$ slab per
coordinate, written atomically or reduced per block. With translational
invariance, three of four centres suffice and the fourth's contribution is
minus the sum. Memory: at 500 AOs and 30 atoms the $\Fsk{x}$ alone are
180 MB, and the CPHF's responses and right-hand sides are small beside
them.

\paragraph{Since written.} The one-electron part is
\code{overlap\_derivative}, \code{kinetic\_derivative} and
\code{potential\_derivative}, the $[n_{\mathrm{atom}}, 3, n_{\mathrm{AO}},
n_{\mathrm{AO}}]$ tensors of \code{src/integrals/one\_electron/gradient};
the two-electron part, over density fitting, is \code{df\_jk\_derivative}
(\code{src/integrals/eri/density\_fitting/gradient}). \code{src/hessian}
assembles $\Fsk{x}$ from them.

%======================================================================
\section{Density fitting}
\label{sec:rhf-df}

The tree's DF-RHF gradient fits $J$ and $K$ with a Coulomb-metric auxiliary
basis, $M_{PQ} = (P|Q) = (LL^{\mathsf T})_{PQ}$ and $T_{\mu\nu,P} =
(\mu\nu|P)$. The DF Hessian has the same structure as sections
\ref{sec:response} and \ref{sec:hessian}, with $G$ replaced by the DF $G$
everywhere (in $F$, in the CPHF product \eqref{eq:matvec}, and in
$G(\Dsk{y})$), and the explicit and skeleton parts differentiated through
the fit. The Coulomb-metric fit is variational in the fit coefficients, so
no fit response appears; only $M^{-1}$'s derivative does.

\subsection{Coulomb}

$E_J = \frac12\,u^{\mathsf T} M^{-1} u$ with $u = T^{\mathsf T}\mathrm{vec}\,D$
and $d = M^{-1}u$, the vector \code{gradient::df\_rhf} already builds. At
fixed $D$,
\begin{align}
  E_J^x &= u^{x\mathsf T} d - \tfrac12\,d^{\mathsf T} M^x d,\\
  E_J^{xy} &= u^{xy\mathsf T} d - \tfrac12\,d^{\mathsf T} M^{xy} d
  + \bigl(u^x - M^x d\bigr)^{\mathsf T} M^{-1} \bigl(u^y - M^y d\bigr),
  \label{eq:dfj}
\end{align}
with $u^x = T^{x\mathsf T}\mathrm{vec}\,D$. The last term is new
relative to the gradient: a $3N\times3N$ Gram matrix of the vectors
$L^{-1}(u^x - M^x d)$, each of length $n_{\mathrm{aux}}$: one triangular
solve per coordinate and one GEMM.

\subsection{Exchange}

$E_x = -\sum_{ij} T_{ij}^{\mathsf T} M^{-1} T_{ij}$ with $T_{ij,P} = (ij|P)$
at fixed MO coefficients and $c_{ij} = M^{-1}T_{ij}$, the fitted occupied
pairs of the gradient. At fixed $C$,
\begin{align}
  E_x^x &= -\sum_{ij}\bigl(2\,T^{x\mathsf T}_{ij} c_{ij}
          - c_{ij}^{\mathsf T} M^x c_{ij}\bigr),\\
  E_x^{xy} &= -\sum_{ij}\Bigl(2\,T^{xy\mathsf T}_{ij} c_{ij}
          - c_{ij}^{\mathsf T} M^{xy} c_{ij}
          + 2\,\bigl(T^x_{ij} - M^x c_{ij}\bigr)^{\mathsf T} M^{-1}
              \bigl(T^y_{ij} - M^y c_{ij}\bigr)\Bigr).
  \label{eq:dfk}
\end{align}
The cross term is the DF Hessian's bottleneck: it needs the half-transformed
derivative three-centre integrals $T^x_{ij,P}$ for every coordinate,
$3N\,n_o^2\,n_{\mathrm{aux}}$ doubles in all. It should be built per atom
(only that atom's AOs and aux functions contribute to $T^x$) and contracted
into the Hessian one pair of atom blocks at a time, never held whole.

These are the fixed-orbital (explicit) parts. Their skeleton analogues,
the derivative of the DF Fock matrix at fixed $D$, follow the same way:
$J^{(x)} = T^x d + T\,M^{-1}(u^x - M^x d)$, and the exchange analogue
with $c_{ij}$.

\subsection{What the fit adds to the kernel list}

$(\mu\nu|P)^{xy}$ and $(P|Q)^{xy}$, the second-derivative three- and
two-centre integrals, contracted with $D\otimes d$, $d\otimes d$ and the
exchange analogues of \code{gradient::df\_rhf}'s $\Gamma^K$ and $W^K$;
and $(\mu\nu|P)^x$ and $(P|Q)^x$ as \emph{uncontracted} per-coordinate
slabs (or contracted over one index only) for \eqref{eq:dfj},
\eqref{eq:dfk} and the skeleton Fock. The existing
\code{three\_center\_gradient} and \code{metric\_gradient} compute the
right integrals and contract them to scalars, like their direct
counterparts. The aux side carries shells two $l$ above the orbital side
(to i), so second derivatives push the aux Hermite expansion to $l = 8$;
that is a strong argument for using translational invariance to move the
derivatives off the aux centre wherever possible.

\paragraph{Since written.} The two-electron explicit part, \eqref{eq:dfj}
and \eqref{eq:dfk} with their cross terms built one atom block at a time,
is \code{hessian::df\_rhf\_two\_electron} (\code{src/hessian}).

%======================================================================
\section{The algorithm}
\label{sec:rhf-algorithm}

\begin{algorithm}[ht]
\caption{RHF analytic Hessian (symmetric form \eqref{eq:hess-sym})}
\label{alg:hessian}
\begin{algorithmic}[1]
\Require converged $C$, $\eps$, $n_o$; $D$, $W$; the basis (and aux basis)
\Ensure $\Exy$, $3N\times3N$
\State $H \gets \Exy_{\mathrm{expl}}$: second-derivative integrals contracted
       with $D$, $-W$, $\frac12 D\otimes D$ \eqref{eq:explicit}; add $V_{nn}^{xy}$
\For{each coordinate $x$ (batched)}
  \State $S^x$, $\Fsk{x}$ as AO matrices \Comment{derivative J/K build}
  \State $S^x_{pq} \gets C^{\mathsf T} S^x C$,\;
         $\Fsk{x}_{pq} \gets C^{\mathsf T}\Fsk{x}C$ (occupied columns only)
  \State $\Dsk{x} \gets 2C_oS^x_{oo}C_o^{\mathsf T}$, then $G(\Dsk{x})$
         \Comment{one J/K build}
  \State $B^x_{ai} \gets -\Fsk{x}_{ai} + \eps_iS^x_{ai} + G(\Dsk{x})_{ai}$
\EndFor
\State solve $MU^x = B^x$ for all $x$ together by preconditioned CG
       \eqref{eq:matvec}, preconditioner $(\eps_a-\eps_i)^{-1}$
\State $H \mathrel{+}= -4\,B^{\mathsf T}U
  - 2(\Sigma + \Sigma^{\mathsf T})
  + 2\,\Omega + \Theta$, where
  $\Sigma_{xy} = \sum_{ij}S^x_{ij}\Fsk{y}_{ij}$,
  $\Omega_{xy} = \sum_{ij}S^x_{ij}S^y_{ij}(\eps_i+\eps_j)$,
  $\Theta_{xy} = \tr[\Dsk{x}G(\Dsk{y})]$
\State check $\|H - H^{\mathsf T}\|$ and the translational sum rule; symmetrize
\end{algorithmic}
\end{algorithm}

\paragraph{Cost.} The J/K builds dominate: $3N$ for the $G(\Dsk{x})$, and
$3N$ per CG iteration, typically 10--20 iterations to $10^{-8}$ in the
residual. That is $3N$ times a short SCF. The derivative J/K build for the
$\Fsk{x}$ costs a small multiple of one gradient. The second-derivative
ERI contraction is one pass over the quartets, more expensive per quartet
than a gradient but done once. For a molecule of 30 atoms the CPHF is
about 90 batched solves, and it is the term to make fast. (These are
expectations from the operation counts, not measurements.)

\paragraph{The CPHF solver.} All $3N$ systems share $M$, so a block CG (or
plain CG run in lockstep, with per-column convergence masks) turns the
$3N$ J/K builds per iteration into one batched build over $3N$ densities.
The integral-direct J/K and the DF J/K both take one density per call
today; a multi-density build is the natural extension, and it serves
TDHF and polarizabilities too. Memory per column is $n_vn_o$ for each of
$U$, $B$, the residual and the search direction: small.
\code{response::zvector} is the tree's closest existing linear solver over
a response product and is the model for its interface and stopping rules.

\paragraph{Since written.} \code{hessian::df\_rhf} is
Algorithm~\ref{alg:hessian} end to end over density fitting
(\code{src/hessian/README.md}, ``The whole Hessian'').

%======================================================================
\section{Checks}
\label{sec:rhf-checks}

\paragraph{Formulas.} The design implemented \eqref{eq:cphf},
\eqref{eq:hess-ao} and \eqref{eq:hess-sym} with PySCF supplying the
integrals ($\Exy_{\mathrm{expl}}$ from \code{partial\_hess\_elec} plus
\code{hess\_nuc}, $\Fsk{x}$ from \code{make\_h1}, $S^x$ from
\code{int1e\_ipovlp}) and solved the CPHF directly with a dense $M$. For
water in 6-31G at a distorted $C_1$ geometry:

\begin{center}
\begin{tabular}{@{}lr@{}}
\toprule
comparison & max.\ abs.\ difference ($E_h/a_0^2$) \\
\midrule
symmetric form \eqref{eq:hess-sym} vs.\ \code{pyscf.hessian.rhf} & $1.9\times10^{-8}$ \\
symmetric form vs.\ central difference of analytic gradients & $7.0\times10^{-8}$ \\
AO form \eqref{eq:hess-ao} vs.\ \code{pyscf.hessian.rhf} & $1.9\times10^{-8}$ \\
asymmetry of the AO form & $3.9\times10^{-10}$ \\
asymmetry of the symmetric form & $1.5\times10^{-15}$ \\
\bottomrule
\end{tabular}
\end{center}

The largest Hessian element is 0.71, so the agreement is at PySCF's own
CPHF tolerance and the finite difference's step error.

\paragraph{For the GPU implementation.} In order of how cheaply they catch
a bug:
\begin{enumerate}
  \item $H = H^{\mathsf T}$ in the AO form, to the CPHF tolerance;
  \item translational invariance: $\sum_B \partial^2E/\partial A_k\partial
        B_l = 0$ for every $A, k, l$. Before symmetrizing
        \eqref{eq:hess-sym}, the sum over the first index $x$ is exact, since
        $\sum_x B^x = 0$ and $\sum_x S^x = 0$ are integral identities; the
        sum over the second index $y$ needs $\sum_y U^y = 0$ and so holds
        only to the CPHF tolerance (unless the last atom's $U$ is taken as
        minus the sum of the others', section~\ref{sec:response}), and
        symmetrizing mixes the two;
  \item rotational invariance at a converged stationary geometry: six
        zero eigenvalues (five for a linear molecule) after mass-weighting;
  \item each term against PySCF separately (\code{partial\_hess\_elec}
        for the explicit part, \code{make\_h1} for $\Fsk{x}$, the
        \code{mo1} of \code{solve\_mo1} for $U$), as the tree's gradient
        tests compare each term; for DF, against \code{pyscf.df.hessian.rhf}
        over the \emph{same} aux set, not against the direct Hessian;
  \item the whole against a central difference of \code{gradient::rhf}
        or \code{gradient::df\_rhf}, which tests the tree against itself.
\end{enumerate}

%======================================================================
\section{What to build, in order}
\label{sec:rhf-plan}

\begin{enumerate}
  \item \textbf{The semi-numerical Hessian.} A central difference of the
        existing analytic gradient: $6N$ SCF + gradient runs at
        $10^{-3}$ to $10^{-4}$ bohr steps, accurate to about $10^{-6}$
        with tight SCF convergence. It needs no new kernel, gives
        frequencies now, and is the reference the analytic one is tested
        against. It belongs beside the geometry-optimization step as a
        pipeline step.
  \item \textbf{Matrix-valued first derivatives.} $S^x$, $h^x$ and the
        derivative J/K as $3N$ AO matrices (section~\ref{sec:skeleton}),
        and the DF analogues. With these alone the CPHF right-hand sides
        exist.
  \item \textbf{The batched CPHF solver} over a multi-density J/K build,
        testable on its own against \code{solve\_mo1}.
  \item \textbf{The response part of $\Exy$} \eqref{eq:hess-sym}: GEMMs and
        dot products once 2 and 3 exist. At this point the Hessian minus
        $\Exy_{\mathrm{expl}}$ can be checked against PySCF.
  \item \textbf{Second-derivative integrals}, one-electron first (cheap,
        easy to verify), then the ERI or DF three- and two-centre
        families, generated by FelizJueves with translational invariance
        built in, and their compile times measured.
  \item \textbf{Frequencies and intensities.} Mass-weighting, projection of
        translations and rotations, harmonic frequencies. IR intensities
        need the dipole derivatives, which reuse the $U^x$:
        $\partial\boldsymbol\mu_{\mathrm{el}}/\partial x =
        -\tr(D\,\boldsymbol\mu^x) - 4\sum_{ai}U^x_{ai}\,\boldsymbol\mu_{ai}
        + 2\sum_{ij} S^x_{ij}\,\boldsymbol\mu_{ij}$ (with $\boldsymbol\mu$
        the dipole integrals $\langle\mu|\mathbf r|\nu\rangle$, electrons
        carrying charge $-1$), plus the nuclear charge $Z_A\,\mathbf e_k$.
\end{enumerate}

Steps 2--4 also serve any other second-order RHF property with a nuclear
perturbation, and the batched CPHF serves electric ones (polarizabilities)
with $B$ replaced by the dipole integrals.

\paragraph{Since written.} Step 1 is \code{src/pipeline/hessian}'s
numerical method; step 6's dipole derivative and IR intensities are
\code{hessian::rhf\_dipole\_derivative} and
\code{hessian::infrared\_intensities}, and its harmonic analysis is
\code{geometry::harmonic\_analysis}.

%======================================================================
\part{CASSCF}

\section{Scope, notation and the energy}
\label{sec:notation}

Orbitals are real, the active space is complete, and the CI vector lives in
a fixed determinant basis. MO indices: core $i,j$; active $t,u,v,w$;
virtual $a,b$; general $p,q,r,s$. AO indices $\mu\nu\lambda\sigma$; aux
indices $P,Q$. Nuclear coordinates are $x$ and $y$, $3N$ of them.

With the 1- and 2-RDMs $D_{pq}$ and $\Gamma_{pqrs}$ over all MOs (core
blocks closed-shell, active blocks from the CI vector), the energy is
\begin{equation}
  E = \sum_{pq} h_{pq} D_{pq}
  + \tfrac12\sum_{pqrs}(pq|rs)\,\Gamma_{pqrs} + V_{nn},
  \label{eq:energy}
\end{equation}
and the generalized Fock matrix, \code{response::fock}'s, is
\begin{equation}
  F_{pq} = \sum_r h_{pr}D_{rq} + \sum_{rst}(pr|st)\,\Gamma_{qrst}.
  \label{eq:fock}
\end{equation}
Only its first $n_{\mathrm{occ}}$ columns are nonzero. The parameters are
$\lambda = (\kappa, c)$: the non-redundant rotations $\kappa$ (core--active,
core--virtual, active--virtual), $C(\kappa) = C\,e^{K}$ with $K$
antisymmetric, and the CI vector. Their gradient is the orbital gradient
$g_{pq} = 2(F_{pq} - F_{qp})$ and the CI gradient $2P(H-E)c$, $P = 1 -
cc^{\mathsf T}$. Both vanish at convergence.

Two energies are in scope:
\begin{itemize}
  \item \textbf{state-specific} (SS): one root, $E$ stationary in all of
        $\lambda$ (section~\ref{sec:ss});
  \item \textbf{state-averaged} (SA): root $I$'s energy $E_I$ at the
        orbitals that make $E_{\mathrm{SA}} = \sum_r w_r E_r$ stationary,
        with every root an eigenvector of the same active-space
        Hamiltonian $H_0$ (section~\ref{sec:sa}).
\end{itemize}

\paragraph{Derivative integrals.} As in Part~I
(section~\ref{sec:rhf-notation}), $h^x$, $(\mu\nu|\lambda\sigma)^x$, $S^x$
and their second derivatives are taken at fixed AO coefficients, and
transformed to the MO basis with the unperturbed $C$: $S^x_{pq} =
C_p^{\mathsf T}S^xC_q$. A quantity in brackets, $F^{[x]}$, $H^{[x]}$,
$g^{[x]}$, means the same expression as $F$, $H$, $g$ evaluated at the
derivative MO integrals $(h^x_{pq}, (pq|rs)^x)$ in place of $(h_{pq},
(pq|rs))$, at the unperturbed RDMs.

%======================================================================
\section{The gradient, as the tree computes it}
\label{sec:gradient}

\code{mcscf::df\_gradient} (\code{src/mcscf/gradient/README.md}) computes
\begin{equation}
  E^x = \sum D^{\mathrm{AO}}_{\mu\nu} h^x_{\mu\nu}
  + \tfrac12\sum \Gamma^{\mathrm{AO}}_{\mu\nu\lambda\sigma}
    (\mu\nu|\lambda\sigma)^x
  - \sum W_{\mu\nu} S^x_{\mu\nu} + V_{nn}^x,
  \qquad W = C\,\sym(F)\,C^{\mathsf T},
  \label{eq:gradient}
\end{equation}
with the connection $C(x) = C(1 - \frac12 S^x_{\mathrm{MO}}\,\delta x)$,
so that $\partial E/\partial U_{pq} = 2F_{pq}$ turns the connection into
$-\sum F_{pq}S^x_{pq}$. For a state average the same terms are taken at
root $I$'s Lagrangian densities: root $I$'s RDMs, plus the transition RDMs
of $z_c$, plus the orbital response of $z_o$
(section~\ref{sec:sa}).

The Hessian is $d/dy$ of \eqref{eq:gradient}, or equivalently the second
derivative of the energy along the parameter-fixed path plus the response.
The second view is the one that organizes the work, and it starts with the
path.

%======================================================================
\section{The connection, to second order}
\label{sec:connection}

At fixed parameters the MOs must stay orthonormal as the basis moves. With
$s(x) = C^{\mathsf T}S(x)C - 1$, the symmetric (L\"owdin) connection is
\begin{equation}
  C(x) = C\,T(x), \qquad
  T = (1+s)^{-1/2} = 1 - \tfrac12 s + \tfrac38 s^2 - \cdots
\end{equation}
Expanding $s = S^x\,\delta x + \frac12 S^{xy}\,\delta x\,\delta y$ (MO
basis throughout) gives the first and second derivatives of $T$,
\begin{equation}
  \tau^x = -\tfrac12 S^x, \qquad
  \tau^{xy} = -\tfrac12 S^{xy} + \tfrac38\bigl(S^xS^y + S^yS^x\bigr).
  \label{eq:tau}
\end{equation}
The tree's gradient uses $\tau^x$ only. \textbf{The linearized connection
$1 - \frac12 S^x\delta x$ is not orthonormal at second order}, and a
Hessian built on it is wrong by terms in $S^xS^y$. Any connection that is
orthonormal to second order gives the same final Hessian at a stationary
point; the pieces differ. The symmetric one is used here because it is the
one the gradient and \code{response::fock} already assume.

\paragraph{The fixed-parameter second derivative.} Write $E(T)$ for the
energy \eqref{eq:energy} with every MO integral transformed by a general,
not necessarily orthogonal, $T$: $h \to T^{\mathsf T}hT$, and $T$ on each
index of $(pq|rs)$. Then $\partial E/\partial T_{pq} = 2F_{pq}$ at $T =
1$, and its second derivative along two generators $A$ and $B$ is the
energy at the \emph{doubly one-index transformed} integrals,
\begin{equation}
  E^{(2)}[A,B] = \sum D_{pq}\bigl(A^{\mathsf T}hB + B^{\mathsf T}hA\bigr)_{pq}
  + \tfrac12\sum\Gamma_{pqrs}\,(pq|rs)^{[A,B]},
\end{equation}
$(pq|rs)^{[A,B]}$ being the sum over every way of applying $A$ to one
index and $B$ to another. This is the machinery behind the orbital
Hessian, \code{response::Hessian}, used here with symmetric generators
instead of antisymmetric ones. Collecting terms, the second derivative of
the energy at fixed $\lambda$ is
\begin{equation}
\boxed{\;
\begin{aligned}
  E_{xy} ={}& \underbrace{\sum D^{\mathrm{AO}} h^{xy}
     + \tfrac12\sum\Gamma^{\mathrm{AO}}(\mu\nu|\lambda\sigma)^{xy}
     + V_{nn}^{xy}}_{\text{explicit}}
   + 2\sum_{pq}F_{pq}\,\tau^{xy}_{pq}\\
  &+ 2\sum_{pq}\bigl(F^{[x]}_{pq}\,\tau^y_{pq} + F^{[y]}_{pq}\,\tau^x_{pq}\bigr)
   + E^{(2)}[\tau^x, \tau^y].
\end{aligned}\;}
\label{eq:exy}
\end{equation}
The second term reduces to $-\tr(WS^{xy})$ plus the $\frac38$ terms, the
RHF $-WS^{xy}$ plus its second-order correction. The third is the
derivative-integral Fock matrix against the first-order connection. The
last is the connection acting twice. For RHF these collapse into the
$S^{x}_{ij}$ terms of the RHF Hessian's symmetric form.

%======================================================================
\section{State-specific CASSCF}
\label{sec:ss}

\subsection{The Hessian}

$E(x, \lambda)$ is the energy at geometry $x$ and parameters $\lambda$
along the connection of section~\ref{sec:connection}, with the orbitals
\begin{equation}
  C(x, \kappa) = C\,T(x)\,e^{K(\kappa)}:
  \label{eq:order}
\end{equation}
the connection first, the rotation second. The order changes the pieces,
not the total. Every per-coordinate formula below, and the $z$-density
shortcut of section~\ref{sec:sa}, assumes this one. At the solution
$E_\lambda = 0$ at every geometry, so
\begin{equation}
  \frac{d E}{dx} = E_x, \qquad
  \frac{d^2 E}{dx\,dy} = E_{xy} + E_{x\lambda}\,\lambda^y,
  \qquad E_{\lambda\lambda}\,\lambda^y = -E_{\lambda y},
\end{equation}
and therefore
\begin{equation}
\boxed{\;
  \Exy = E_{xy} - E_{x\lambda}\,E_{\lambda\lambda}^{-1}\,E_{\lambda y}
       = E_{xy} - \lambda^{x\mathsf T}E_{\lambda\lambda}\,\lambda^y.\;}
  \label{eq:ss}
\end{equation}
The response term is negative semi-definite when $E_{\lambda\lambda}$ is
positive definite, a minimum in the parameters, as for RHF. A
state-specific \emph{excited} root is a saddle point in $\lambda$: the
response term then has no sign, and a conjugate-gradient solve of
$E_{\lambda\lambda}\lambda^y = -E_{\lambda y}$ is not guaranteed to
converge.
$E_{xy}$ is \eqref{eq:exy}. $E_{\lambda\lambda}$ is the CASSCF parameter
Hessian (orbital--orbital, orbital--CI, CI--CI), the one the second-order
optimizer and the $Z$-vector already apply. $E_{x\lambda}$ is the
right-hand side below. The design checked \eqref{eq:ss} numerically
against a central difference of PySCF's analytic CASSCF gradient.

\subsection{The CP-CASSCF right-hand side}

$E_{\lambda y}$ is the parameter gradient differentiated at fixed
parameters along the connection:
\begin{align}
  \text{orbital:}\quad & g^{[y]}_{pq} + \bigl(\text{$g$ at the integrals
     one-index transformed by } \tau^y\bigr)_{pq},
     \qquad g^{[y]}_{pq} = 2\bigl(F^{[y]}_{pq} - F^{[y]}_{qp}\bigr),
     \label{eq:rhs-orb}\\
  \text{CI:}\quad & 2P\bigl(\tilde H^{y} - \langle c|\tilde H^{y}|c\rangle\bigr)c,
     \qquad \tilde H^{y} = H^{[y]} + H[\tau^y],
     \label{eq:rhs-ci}
\end{align}
where $H[\tau^y]$ is the active-space Hamiltonian at the $\tau^y$
one-index transformed integrals (core included, through $h_{\mathrm{eff}}$).
$F^{[y]}$ needs the derivative integrals with all four indices in the
occupied and general spaces, contracted with $D$ and $\Gamma$: a
\emph{derivative J/K build} over the core density and the active density,
plus the active-pair derivative integrals $(pt|uv)^y$. These are the
CASSCF version of the RHF skeleton Fock matrices
(section~\ref{sec:skeleton}), and like those they are matrices per
coordinate, not the scalars the tree's derivative kernels return.

\subsection{Solving it: the \texorpdfstring{$Z$}{Z}-vector's operator, \texorpdfstring{$3N$}{3N} times}
\label{sec:3n}

$E_{\lambda\lambda}\lambda^y = -E_{\lambda y}$ is the $Z$-vector equation
with a different right-hand side. \code{response::zvector}
(\code{src/response/zvector/README.md}) already solves it with a nonzero CI
part, by the Schur complement on the CI-relaxed orbital Hessian $\tilde H$
of \code{response::RelaxedHessian}:
\begin{equation}
  \tilde H\,\kappa^y = -g_y - H_{oc}w_y,
  \qquad w_y = -(H_0 - E)^{-1}P\,y_y,
  \qquad c^y = -(H_0 - E)^{-1}P\,H'(\kappa^y)\,c + w_y,
  \label{eq:schur}
\end{equation}
with $(g_y, y_y)$ the orbital and CI parts of $E_{\lambda y}$. The scaling
trap that README documents applies unchanged: its CI rows are
\emph{unweighted and unscaled}, so a right-hand side written with the
$2P(H-E)c$ convention of \eqref{eq:rhs-ci} enters halved. Assemble the
response term as $\Lambda^{\mathsf T}E_{\lambda\cdot}$, the solutions
against the right-hand sides, not as $-\Lambda^{\mathsf T}E_{\lambda\lambda}
\Lambda$. The operator's scaling then cancels between the solve and its
right-hand side and never reaches the Hessian.

What changes is the count. One $\tilde H$ product is J/K builds plus an
inner conjugate-gradient CI solve (up to 200 sigmas per root), so $3N$
right-hand sides through the Schur complement nest one iterative solve
inside another, $3N$ times. Two remedies, to be measured rather than
assumed:
\begin{enumerate}
  \item \textbf{Batch the Schur complement.} All $3N$ share $\tilde H$, so
        block CG applies one product to $3N$ vectors: one batched J/K
        build and $3N$ CI solves per iteration.
  \item \textbf{Solve the coupled system unreduced.} PCG on the full
        $(\kappa, c)$ vector with the block-diagonal preconditioner
        $(\operatorname{diag}H_{oo},\,(H_0 - E_0)_{\mathrm{diag}})$, as PySCF
        solves its CP-MCSCF equations. No inner solve; each iteration is
        one orbital Hessian product, one $H'(\kappa)c$ and one sigma per
        right-hand side.
\end{enumerate}

\paragraph{What the J/K builds see.} Wang et al.\ count two J and two K
builds per PCG step of TeraChem's CP-MCSCF, on two densities built from
the trial vector $p$: the core--virtual density $D_{CV}$ (only the
$p_{\mathrm{core,virt}}$ block) and the generalized active density
\[
  D_{\mathrm{Gen}} =
  \begin{pmatrix}
    0 & p_{\mathrm{core,act}}(2 - \gamma) & 2p_{\mathrm{core,virt}}\\
    0 & 0 & \gamma\,p_{\mathrm{act,virt}}\\
    0 & 0 & 0
  \end{pmatrix}
\]
(their Eqs.~14 and 15). Only the symmetric parts contribute, but the
symmetrized densities are not positive semi-definite, so they cannot be
written $XX^{\mathsf T}$ for DF-K-MO. Their fix is an SVD,
$\tilde D = U\Sigma V^{\mathsf T}$, giving left and right pseudo-coefficients
$C_L = U\Sigma^{1/2}$ and $C_R = V\Sigma^{1/2}$, and an exchange build
$K(C_L, C_R)$ over the fitted pairs. \textbf{The symmetrized $\tilde D$ has
rank up to $2n_{\mathrm{occ}}$, not $n_{\mathrm{occ}}$.} The unsymmetrized
$D_{\mathrm{Gen}}$ and $D_{CV}$ have nonzero rows only in the occupied
block, so their rank is at most $n_{\mathrm{occ}}$ (or $n_{\mathrm{core}}$).
They also come factored, with no SVD: in the AO basis $D = C_o\,(D_{o\cdot}
C^{\mathsf T})$, so $C_L = C_o$ and $C_R = C\,D_{o\cdot}^{\mathsf T}$.
$K(C_L, C_R)$ followed by $K \to \frac12(K + K^{\mathsf T})$ is then
$K(\tilde D)$ at half the rank. The SVD is a compression to try, not a
requirement. Every density the Hessian adds is of this kind: the $3N$ response
densities per CP-CASSCF iteration, the CI responses' transition RDMs, and
in the state average the $Z$-weighted ones. A batched, left/right DF-K over
SVD factors is the J/K the Hessian wants.

\paragraph{Since written.} \code{response::zvectors} takes the $3N$
right-hand sides as a block, and solves them column by column, by the
batched Schur complement, or by the coupled system unreduced
(\code{src/response/zvector/README.md}, ``Many right-hand sides'').

%======================================================================
\section{State-averaged CASSCF}
\label{sec:sa}

\subsection{The Lagrangian, and why its second derivative needs only \texorpdfstring{$Z$}{Z}}
\label{sec:sa-lagrangian}

Root $I$'s energy is not stationary in the SA parameters. The gradient is
that of
\begin{equation}
  \Lag(x, \lambda, z) = E_I(x, \lambda) + z\cdot g_{\mathrm{SA}}(x, \lambda),
  \qquad g_{\mathrm{SA}} = \partial E_{\mathrm{SA}}/\partial\lambda,
\end{equation}
with $A z = -\partial E_I/\partial\lambda$, $A =
\partial g_{\mathrm{SA}}/\partial\lambda$ the SA Hessian: the tree's
$Z$-vector. $\Lag$ is stationary in $\lambda$ and $z$, and the same
elimination as the NEVPT2 Hessian's $2n+1$ form
(\code{docs/nevpt2/nevpt2.tex}, Eq.~(hess)) gives
\begin{equation}
\boxed{\;
  \Exy_I = \Lag_{xy} + \Lag_{x\lambda}\lambda^y + \lambda^{x\mathsf T}\Lag_{\lambda y}
     + \lambda^{x\mathsf T}\Lag_{\lambda\lambda}\lambda^y,
  \qquad A\,\lambda^x = -\partial_x g_{\mathrm{SA}}.\;}
  \label{eq:sa}
\end{equation}
\begin{itemize}
  \item $\lambda^x$ is the response of the \emph{SA} parameters: $3N$
        solves with the SA Hessian $A$, the same operator as the
        $Z$-vector, against the SA right-hand sides
        \eqref{eq:rhs-orb}--\eqref{eq:rhs-ci} at the averaged densities.
  \item $z$ enters at zeroth order, the gradient's own. No $z^x$, which
        would be $3N$ more solves.
  \item $\Lag_{xy} = E_{I,xy} + z\cdot\partial_{xy}g_{\mathrm{SA}}$: the
        fixed-parameter second derivative \eqref{eq:exy} at root $I$'s
        densities plus the same at $z$'s, the transition RDMs of $z_c$
        and the $z_o$-rotated densities. These are exactly the densities
        the gradient already forms.
  \item $\Lag_{x\lambda} = \partial_x\partial_\lambda E_I
        + z\cdot\partial_x\partial_\lambda g_{\mathrm{SA}}$: root $I$'s
        CP right-hand side, plus the SA Hessian at derivative integrals
        (and connection) contracted with $z$.
  \item $\Lag_{\lambda\lambda} = \partial_{\lambda\lambda}E_I
        + z\cdot\partial_{\lambda\lambda}g_{\mathrm{SA}}$: root $I$'s
        parameter Hessian, plus a \textbf{third derivative of
        $E_{\mathrm{SA}}$ contracted once with $z$}. Only its action
        $\lambda^{x\mathsf T}(\cdot)\lambda^y$ is needed.
\end{itemize}

The third-derivative term is the new object, and the warning the NEVPT2
Hessian carries applies to it verbatim: \textbf{differencing a
Hessian-vector product does not give it} on any pair with an orbital part,
because the operator is the Jacobian of $g$ only at a stationary point. An
analytic version differentiates \code{response::fock}'s gradient twice: J/K
builds at doubly rotated densities, plus the CI legs as transition
densities.

\paragraph{Since written.} The third-derivative action is
\code{mcscf::third\_derivative} (\code{src/mcscf/third\_derivative}), the
covector $a \mapsto T(a, v, z)$ for each vector $v$ of a block.

\subsection{The rotations among the averaged roots}
\label{sec:roots}

The SA energy is invariant to rotations among its roots when their weights
are equal. Root $I$'s is not. They are fixed by the condition that $H_0$
stays diagonal in the roots' span, $\langle c_r|H|c_s\rangle = 0$ for
$r \neq s$. With $c_r \to c_r + \sum_s\theta_{sr}c_s$, $\theta$
antisymmetric, its first derivative gives
\begin{equation}
  \theta^x_{sr} = \frac{\mathcal H^{x}_{rs}}{E_r - E_s},
  \qquad
  \mathcal H^{x}_{rs} = \langle c_r|\tilde H^{x} + H'(\kappa^x)|c_s\rangle,
  \label{eq:theta}
\end{equation}
the orbital-relaxed coupling between the two roots, which is the
derivative-coupling numerator. The CI response outside the span does not
enter it, because $H_0 c_s = E_s c_s$ is orthogonal to that response.

\paragraph{In the gradient.} At equal weights the rotations are redundant in
$E_{\mathrm{SA}}$, and $E_I$ is stationary in them because the roots are
eigenvectors. Their multipliers vanish, and the tree projects them away
(\code{src/response/zvector/README.md}). \textbf{At unequal weights
they are not redundant.} $\partial^2E_{\mathrm{SA}}/\partial\theta_{rs}^2 =
2(w_r - w_s)(E_s - E_r)$, and they couple to the orbitals through
$(w_r - w_s)$ times the $r$--$s$ transition orbital gradient, so they
belong in $A$ and can carry multipliers. A PySCF model of the state average
finds:
\begin{itemize}
  \item \textbf{two roots, any weights}: the multiplier of the one rotation
        is zero to $3\times10^{-7}$ at weights $(0.7, 0.3)$, $(0.9, 0.1)$
        and $(0.35, 0.65)$, against a coupling up to $0.16$. The projected
        gradient is exact. This is observed, not proven;
  \item \textbf{three roots, weights $(0.5, 0.3, 0.2)$}: the rotations
        $\theta_{01}$ and $\theta_{02}$ again carry no multiplier, but
        $\theta_{12}$ carries 0.07--0.15. The projected gradient is
        \textbf{wrong by up to $3.3\times10^{-3}$}, 7\% of root~2's
        gradient. With the rotations in $A$ it is exact to $2\times10^{-6}$.
\end{itemize}
So the gradient as it then was is right at equal weights and for two
roots, and wrong for three or more roots at unequal weights. That case is
outside PySCF's analytic SA gradient too, which refuses unequal weights.
Until the tree carries the rotations for pairs of unequal weight, it should
refuse that case. Partly equal weights, where some rotations are redundant
and some are not, were not checked.

\paragraph{In the Hessian.} At equal weights the rotations are not in
$\lambda$, and they add a term of their own to \eqref{eq:sa}. $E_I$ is
stationary in $\theta_{sI}$, with $\partial^2E_I/\partial\theta_{sI}^2 =
2(E_s - E_I)$ and total derivative $2\mathcal H^x_{Is}$ along the response.
Eliminating $\theta$ then gives
\begin{equation}
\boxed{\;
  \Exy_I \mathrel{+}= \sum_{s\neq I}
    \frac{\mathcal H^x_{Is}\mathcal H^y_{sI} + \mathcal H^y_{Is}\mathcal H^x_{sI}}
         {E_I - E_s}
  = 2\sum_{s\neq I}\frac{\mathcal H^x_{Is}\mathcal H^y_{Is}}{E_I - E_s}.\;}
  \label{eq:rot}
\end{equation}
\textbf{The factor 2 is not optional.} With only one copy of the product,
the Hessian is wrong by half this term, 0.02 to 0.24 in the design's
checks. With three equal roots the term is as large as the Hessian itself.
At unequal weights the rotations are part of $\lambda$, their response is
inside $\lambda^x$, and \eqref{eq:sa} is complete without \eqref{eq:rot}.

Either way the term \textbf{diverges as $E_s \to E_I$}. That is the
physics: at a conical intersection the adiabatic energy is not
differentiable and has no Hessian. \code{RelaxedHessian} already drops the
other-root part of the CI response for a gap under $10^{-8}$ hartree. A
Hessian must instead refuse, with the gap in the message, below a threshold
on $|\mathcal H^x_{Is}|/|E_I - E_s|$ that is far larger. The tree's weights
are recovered from a solution and need not be equal, so both paths are
needed.

\paragraph{Since written.} Every pair of unequal weight is a parameter of
$A$ with a multiplier of its own (\code{src/response/zvector/README.md},
``The rotations among the roots''), and the equal-weight term
\eqref{eq:rot}, with the rotations \eqref{eq:theta}, is
\code{mcscf::root\_rotation} (\code{src/mcscf/root\_rotation}).

%======================================================================
\section{The explicit part over density fitting}
\label{sec:df}

The explicit part of \eqref{eq:exy} is the second-derivative integrals
contracted with $D^{\mathrm{AO}}$ and $\Gamma^{\mathrm{AO}}$. One-electron:
as RHF, at the CASSCF $D$. Two-electron, over a Coulomb-metric fit with
$T_{a,P} = (a|P)$ for an AO pair $a = (\mu\nu)$, $N = M^{-1}$ and fitted
pairs $c = TN$: for any symmetric pair-space density $\Gamma_{ab}$ (the
energy being $E_2 = \frac12\sum_{ab}\Gamma_{ab}(TNT^{\mathsf T})_{ab}$),
\begin{equation}
\boxed{\;
  E_2^{xy} = \sum_{aP}\tilde\Gamma_{aP}\,T^{xy}_{aP}
  - \tfrac12\sum_{PQ}\hat W_{PQ}\,M^{xy}_{PQ}
  + \sum_{ab}\Gamma_{ab}\bigl[R^x N R^{y\mathsf T}\bigr]_{ab},
  \qquad R^x = T^x - c\,M^x,\;}
  \label{eq:df2}
\end{equation}
with $\tilde\Gamma = \Gamma c$ and $\hat W = c^{\mathsf T}\Gamma c$, the
same intermediates as the gradient. For $\Gamma = D\otimes D$ this is
Part~I's DF Coulomb second derivative \eqref{eq:dfj}. The first two terms
are the gradient's contractions with second-derivative integrals. The third
is new, and it is affordable only because $\Gamma$ factorizes. The tree's
gradient already splits $E_2$ into the pieces that make it so
(\code{src/mcscf/gradient/README.md}):
\begin{itemize}
  \item \textbf{J}, core--core and core--active: $\Gamma = D\otimes D -
        D_a\otimes D_a$. The cross term is two Gram matrices of vectors
        $u^x - M^x d$, length $n_{\mathrm{aux}}$, as for RHF.
  \item \textbf{K}, core with core and with active: $-\sum_i\sum_{pq}
        G_{pq}(ip|iq)$. The cross term needs $R^x$ half-transformed to
        $(i, p)$ pairs, $3N\,n_{\mathrm{core}}n_{\mathrm{occ}}n_{\mathrm{aux}}$
        in all. Build it per atom.
  \item \textbf{active--active}: $\frac12\sum\Gamma_{tuvw}(tu|vw)$. The
        cross term is $\sum\Gamma_{tuvw}(R^x_{tu})^{\mathsf T}N R^y_{vw}$
        with $R^x$ over active pairs, $3N\,n_{\mathrm{act}}^2n_{\mathrm{aux}}$:
        small, and the cleanest of the three.
\end{itemize}
The left/right forms of Wang et al.'s SI (their Eqs.~S9 and S15, the DF-J
and DF-K gradients at independent $D_L$ and $D_R$) are the gradients these
cross terms differentiate. Their DF-K-MO variant (Eqs.~S16--S19), which
contracts over MO pairs and never forms an $n_{\mathrm{aux}}n_{\mathrm{AO}}^2$
transition tensor, is the right shape for the $R^x$ half-transforms too.

The connection terms of \eqref{eq:exy} need $F^{[x]}$ (the derivative J/K
builds of section~\ref{sec:ss}) and $E^{(2)}[\tau^x,\tau^y]$, which is the
orbital-Hessian contraction with symmetric generators. Where $\tau^x$ and
$\tau^y$ fall on different pairs of a core density, it is J/K builds on
$C\tau^x$-rotated densities, $3N$ of them, contracted pairwise. Where both
fall on the same pair, the core part is $\tau^{x\mathsf T}J(D)\tau^y$-type
contractions of the unperturbed J/K. Where they meet the active 2-RDM, it
needs the $(pq|tu)$- and $(pt|qu)$-type integrals with two general
indices: the \code{ppaa} and \code{papa} blocks that \code{response::Hessian}
already contracts for antisymmetric generators.

\paragraph{Since written.} \eqref{eq:df2}, split as above, is
\code{hessian::df\_casscf\_two\_electron}, and \eqref{eq:exy} built on it
is \code{hessian::df\_casscf\_exy}; $F^{[x]}$ and the right-hand sides
\eqref{eq:rhs-orb}--\eqref{eq:rhs-ci} are \code{mcscf::df\_derivatives}
(\code{src/mcscf/derivative}).

%======================================================================
\section{The algorithm}
\label{sec:algorithm}

\begin{algorithm}[ht]
\caption{State-specific CASSCF analytic Hessian, Eq.~\eqref{eq:ss}}
\label{alg:ss}
\begin{algorithmic}[1]
\Require converged $C$, $c$; $D$, $\Gamma$, $F$; basis and aux basis
\Ensure $\Exy$, $3N\times3N$
\State $H \gets$ explicit part: second-derivative integrals against
       $D^{\mathrm{AO}}$, $\Gamma$ split as in section~\ref{sec:df};
       $+V_{nn}^{xy}$
\For{each coordinate $x$ (batched)}
  \State $S^x$, $h^x$, derivative J/K at the core and active densities,
         $(pt|uv)^x$ \Comment{matrix-valued derivatives}
  \State $\tau^x \gets -\frac12 S^x$;\; $F^{[x]}$;\;
         $E_{\lambda x}$ by \eqref{eq:rhs-orb}--\eqref{eq:rhs-ci}
\EndFor
\State $H \mathrel{+}= 2F\cdot\tau^{xy} + 2(F^{[x]}\cdot\tau^y + F^{[y]}\cdot\tau^x)
       + E^{(2)}[\tau^x,\tau^y]$ \Comment{connection, Eq.~\eqref{eq:exy}}
\State solve $E_{\lambda\lambda}\lambda^x = -E_{\lambda x}$ for all $x$:
       block PCG, batched products \Comment{the $Z$-vector operator}
\State $H \mathrel{+}= -\Lambda^{\mathsf T}E_{\lambda\lambda}\Lambda$
       \;$=\; \Lambda^{\mathsf T}E_{\lambda\cdot}$, \;$\Lambda = (\lambda^1\cdots\lambda^{3N})$
\State check symmetry and the translational sum rule; symmetrize
\end{algorithmic}
\end{algorithm}

The state average adds, in order of difficulty: the $z$-weighted copies
of lines 1--4 (the gradient's densities, one more pass); the root-rotation
responses \eqref{eq:theta}, closed form once $\lambda^x$ is known; the
$z$-contracted SA Hessian at derivative integrals in $\Lag_{x\lambda}$;
and the third-derivative action $\lambda^{x\mathsf T}(z\cdot
\partial_{\lambda\lambda}g_{\mathrm{SA}})\lambda^y$, $3N(3N+1)/2$ pairs,
each a doubly-rotated J/K build unless it is organized as $3N$ actions
on $\lambda^y$ followed by dot products.

\paragraph{Cost.} As for RHF, the response dominates: $3N$ CP-CASSCF
solves, each costing about one $Z$-vector, against a gradient that costs
one. A CASSCF Hessian is therefore about $3N$ CASSCF gradients plus the
explicit part, and the semi-numerical route costs $6N$ CASSCF runs
\emph{plus} $6N$ gradients. The analytic route wins by more than the
factor of two that counting suggests, because each displaced CASSCF
re-converges the orbitals from scratch. That is an expectation from the
operation counts, not a measurement.

\paragraph{Since written.} \code{hessian::df\_casscf} is
Algorithm~\ref{alg:ss}, with its CP-CASSCF step on its own as
\code{hessian::df\_casscf\_cp}; the state average is
\code{hessian::df\_sacasscf}, its $z$ terms \code{hessian::df\_sacasscf\_z}
(\code{src/hessian/README.md}).

%======================================================================
\section{Differentiability}
\label{sec:diff}

Every condition a smooth CASSCF gradient needs, the Hessian needs too,
and more sharply, because the response now enters quadratically:
\begin{itemize}
  \item a non-singular $E_{\lambda\lambda}$ (or $A$), a true minimum in the
        orbital--CI space. A near-singular one, such as an active orbital
        with occupation near 0 or 2 that is nearly redundant with the core
        or virtual space, makes $\lambda^x$ large and the Hessian
        ill-conditioned before it makes the gradient wrong;
  \item no root flipping in the state average and no active-space swap
        along the displacements. The semi-numerical Hessian is exposed to
        this, the analytic one is not, and the disagreement between them
        is a diagnostic;
  \item for the state average, a gap to every other averaged root
        (section~\ref{sec:roots}).
\end{itemize}

%======================================================================
\section{Checks}
\label{sec:checks}

\paragraph{The state-specific formula.} The design's check evaluates the
CASSCF energy $E(x,\lambda)$ directly: L\"owdin connection to all orders,
orbitals $C(x)e^{K(\kappa)}$, and the CI vector $c_0 + Q\theta$
normalized, with $Q$ an orthonormal basis of $c_0$'s complement. It takes
$E_{xy}$, $E_{x\lambda}$ and $E_{\lambda\lambda}$ by central finite
differences at fixed parameters and assembles \eqref{eq:ss}. The reference
is a central difference of PySCF's analytic CASSCF gradient, re-converged
at each displaced geometry. For water, CAS(4,4)/6-31G, at a $C_1$ geometry
with both O--H bonds stretched (1.40 and 1.51~\AA), so that the active
space is the two $\sigma/\sigma^*$ pairs with natural occupations near
1.88, 1.84, 0.16 and 0.12:

\begin{center}
\begin{tabular}{@{}lr@{}}
\toprule
comparison & max.\ abs.\ difference ($E_h/a_0^2$) \\
\midrule
formula \eqref{eq:ss} vs.\ FD of PySCF's analytic gradient & $2.8\times10^{-4}$ \\
asymmetry of that FD reference (its own noise) & $2.3\times10^{-4}$ \\
diagonal of \eqref{eq:ss} vs.\ FD of self-relaxed energies & $2.5\times10^{-6}$ \\
$E_{xy}$ alone (no response) vs.\ FD of PySCF's gradient & $4.8\times10^{-1}$ \\
asymmetry of \eqref{eq:ss} & $2\times10^{-16}$ \\
\bottomrule
\end{tabular}
\end{center}

\noindent The largest Hessian element is 0.06, and the response term is
not small: dropping it misses by 0.48. The first row sits at the
reference's own noise, set by PySCF's reachable gradient convergence
($10^{-6}$). The third row replaces PySCF with a reference that has no
solver noise: $E(x,\lambda)$ relaxed to $10^{-9}$ in its gradient at each
displaced geometry by Newton steps, the relaxed energies then differenced.
This checks the structure \eqref{eq:ss} and the connection, but
not the analytic expansions \eqref{eq:exy}--\eqref{eq:rhs-ci} term by
term.

\paragraph{The analytic terms.} A second check takes the derivative AO
integrals (first and second, at fixed AO coefficients) by finite
difference. It assembles $E_{xy}$ by \eqref{eq:exy} and $E_{\lambda x}$ by
\eqref{eq:rhs-orb}--\eqref{eq:rhs-ci} exactly as written, with $E^{(2)}$
over every pair of index positions. It compares them with the
fixed-parameter finite differences above, at the same molecule:

\begin{center}
\begin{tabular}{@{}lrr@{}}
\toprule
term & max.\ size & max.\ abs.\ difference \\
\midrule
$E_{xy}$, \eqref{eq:exy} & 0.44 & $2.5\times10^{-7}$ \\
\quad the same without the $\frac38(S^xS^y + S^yS^x)$ of $\tau^{xy}$ & & $9.3\times10^{-1}$ \\
orbital right-hand side, \eqref{eq:rhs-orb} & 2.0 & $1.5\times10^{-6}$ \\
CI right-hand side, \eqref{eq:rhs-ci} & 0.21 & $2.5\times10^{-7}$ \\
\bottomrule
\end{tabular}
\end{center}

\noindent The second-order connection is not a small correction: without
its $S^xS^y$ terms $E_{xy}$ is wrong by twice its own size.

\paragraph{The state average.} A third check does the same for
SA-CASSCF(4,4) over two or three singlet roots, at the same geometry. The
roots are $(Ve^{A})_{\cdot r}$ over an orthonormal basis $V$ of the singlet
CI space, $A$ holding the root--complement rotations and, as separate
parameters, the root--root rotations $\theta$. The reference solution is
polished by Newton steps to $|g_{\mathrm{SA}}| < 10^{-10}$. The reference
energies are re-relaxed at each displaced geometry and the roots
re-diagonalized in their span. Every piece of \eqref{eq:sa} is a finite
difference of the energy matrix $\langle c_r|H|c_s\rangle$, with the $z$
terms as a directional difference of $E_{\mathrm{SA}}$ along $z$. The
Hessian is compared on its diagonal and along three random directions.
The gradient compares the Lagrangian with the rotations projected (the tree,
PySCF) and with them in $A$ against the differenced relaxed energies. Each
row is the worst over the roots:

\begin{center}
\small
\begin{tabular}{@{}lrrr@{}}
\toprule
weights & gradient, projected & gradient, rotations in $A$ & Hessian \\
\midrule
$(0.5, 0.5)$ & $6\times10^{-7}$ & --- & $3\times10^{-5}$ \\
$(0.7, 0.3)$ & $6\times10^{-7}$ & $6\times10^{-7}$ & $5\times10^{-5}$ \\
$(\frac13, \frac13, \frac13)$ & $2\times10^{-6}$ & --- & $2\times10^{-5}$ \\
$(0.5, 0.3, 0.2)$ & $\mathbf{3.3\times10^{-3}}$ & $2\times10^{-6}$ & $3\times10^{-5}$ \\
\bottomrule
\end{tabular}
\end{center}

\noindent The Hessian column is \eqref{eq:sa}, plus \eqref{eq:rot} at
equal weights. The largest Hessian elements are 0.05 to 0.48 and the
largest gradients 0.04 to 0.16. Without \eqref{eq:rot} the equal-weight
Hessians miss by 0.04 (two roots) and 0.48 (three). With one copy of the
product instead of two they miss by half that. The residuals are
round-off of the nested differences, not a missing term: they grow about
tenfold for each halving of the step ($2\times10^{-3}$ here).
Section~\ref{sec:roots} has what the gradient rows mean.

$\Lag_{x\lambda}$ is also checked analytically in $x$, both parts against
the mixed $(x,\lambda)$ difference, at every weight set above: root $I$'s
right-hand side \eqref{eq:rhs-orb}--\eqref{eq:rhs-ci} at its own densities
(the root--root rotations included at unequal weights) to $6\times10^{-6}$
of a maximum of 1.9, and $z$ times the SA Hessian at the derivative
integrals plus the $\tau^x$ one-index transform to $5\times10^{-6}$ of a
maximum of 0.14. The SA Hessian product in the second is still a difference
in $\lambda$; it is the operator the tree already has. Not checked
analytically: the $z$ part of $\Lag_{xy}$ (\eqref{eq:exy} at $z$'s
densities), the third derivative, and partly equal weights.

\paragraph{Density fitting.} A fourth check repeats the state-specific
checks over a Coulomb-metric fit (cc-pVDZ-JKFIT), at the same molecule.
The DF first-derivative four-index integrals are formed analytically,
$T^xc^{\mathsf T} + cT^{x\mathsf T} - cM^xc^{\mathsf T}$, from
finite-difference $T^x$ and $M^x$. PySCF's DF-CASSCF is only the starting
point. The solution is polished on the check's own DF energy, and the
reference is that energy relaxed by Newton steps on its analytic gradient
to $10^{-10}$:

\begin{center}
\begin{tabular}{@{}lr@{}}
\toprule
comparison & max.\ abs.\ difference \\
\midrule
\eqref{eq:df2} vs.\ FD of $\frac12\Gamma\cdot(TNT^{\mathsf T})$, fixed $\Gamma^{\mathrm{AO}}$ & $3.4\times10^{-8}$ \\
\quad the same without the $R^xNR^{y\mathsf T}$ term & $3.2\times10^{-2}$ \\
J + K + active cross terms (section~\ref{sec:df}) vs.\ generic $\Gamma_{ab}$ & $5\times10^{-14}$ \\
$E_{xy}$, \eqref{eq:exy}, over DF integrals & $2.1\times10^{-7}$ \\
orbital and CI right-hand sides & $1.5\times10^{-6}$, $2.4\times10^{-7}$ \\
\eqref{eq:ss} vs.\ relaxed DF energies (diagonal, 3 directions) & $1.0\times10^{-6}$ \\
\bottomrule
\end{tabular}
\end{center}

\noindent The metric is ill-conditioned, $\kappa(M)\approx2\times10^{6}$
here. An explicit $M^{-1}$ puts $10^{-10}$ of round-off on every energy,
which a second difference at step $10^{-3}$ turns into $10^{-4}$: the first
run of the check missed by that much. Contract through the Cholesky factor
of $M$ in anything that is differenced, the semi-numerical Hessian of
section~\ref{sec:plan} included. The analytic Hessian never differences an
energy and is not affected.

\paragraph{For the GPU implementation.}
\begin{enumerate}
  \item symmetry, and the translational sum rule;
  \item $E_{xy}$ alone against a finite difference of the energy at fixed
        parameters, exactly as the check does: it isolates the
        connection, the bug most likely to survive everything else;
  \item each $\lambda^x$ against the difference of converged parameters
        at displaced geometries (after aligning phases and the connection);
  \item the whole against a central difference of \code{mcscf::df\_gradient},
        SS and each SA root. Include an equal-weight case with three roots,
        where \eqref{eq:rot} is as large as the Hessian, and an unequal-weight
        case;
  \item against OpenMolcas, which computes analytic CASSCF Hessians, over
        the same auxiliary basis for the DF form.
\end{enumerate}

\paragraph{Since written.} \code{tools/emit/casscf\_hessian\_reference.py}
ports these checks onto density fitting throughout, and its data under
\code{test/gtest/data/casscf\_hessian/} reproduce the tables above.

%======================================================================
\section{What to build, in order}
\label{sec:plan}

\begin{enumerate}
  \item \textbf{The semi-numerical Hessian} over \code{mcscf::df\_gradient},
        SS and SA, with the previous geometry's orbitals and CI as guess
        and a check that the active orbitals and the root order did not
        change. It is the reference everything below is tested against,
        and it is only as good as its displaced solves. Building this
        document's check hit two failures that raised no error and gave
        differenced energies and gradients wrong by an order of magnitude.
        A weakly correlated CAS(4,4) near equilibrium never converged.
        Then a displaced CASSCF started from the reference orbitals, which
        are not orthonormal in the displaced overlap metric, converged to an
        energy $O(\delta x)$ off: PySCF does not re-orthonormalize a guess.
        Orthonormalize the guess with the connection, $C(S_{\mathrm{MO}})^{-1/2}$.
        Refuse an unconverged
        displacement.
  \item \textbf{The RHF Hessian's kernels} (Part~I,
        section~\ref{sec:rhf-plan}, items 2, 3 and 5): matrix-valued first
        derivatives, the multi-density J/K, second-derivative integrals.
        The CASSCF Hessian needs all of them, at more densities.
  \item \textbf{The left/right SVD-factored DF-K}, after Wang et al., for
        the non-PSD response and transition densities.
  \item \textbf{The second-order connection} and $E_{xy}$ \eqref{eq:exy},
        checked against fixed-parameter finite differences.
  \item \textbf{The batched CP-CASSCF solve}, the $Z$-vector's operator on
        $3N$ right-hand sides. Measure the batched Schur complement
        against the unreduced coupled PCG before choosing.
  \item \textbf{The SS Hessian} \eqref{eq:ss}. Frequencies of a
        single-state CASSCF.
  \item \textbf{The SA Hessian} \eqref{eq:sa}: the $z$ terms, the root
        rotations (\eqref{eq:rot} at equal weights, in $\lambda$ otherwise)
        with their refusal, and last the third derivative. The NEVPT2
        Hessian needs the same third derivative, so it is built once for
        both. First, and independently of the Hessian, make
        \code{mcscf::df\_gradient} refuse three or more roots at unequal
        weights, or carry their rotations in the $Z$-vector
        (section~\ref{sec:roots}).
\end{enumerate}

\paragraph{Since written.} Item 1 is \code{src/pipeline/hessian}'s
numerical method over the CASSCF chain, which orthonormalizes its guess in the
displaced metric and refuses a displaced solve that left the reference's
solution (\code{src/pipeline/hessian/README.md});
items 6 and 7 are the modules of section~\ref{sec:algorithm}'s ``Since
written''; and \code{mcscf::df\_gradient} carries the unequal-weight
rotations.

%======================================================================
\section*{References}
\begin{itemize}
  \item J.~A. Pople, R. Krishnan, H.~B. Schlegel, J.~S. Binkley,
        \emph{Int.\ J.\ Quantum Chem.\ Symp.}\ \textbf{13}, 225 (1979):
        the RHF second derivatives and CPHF.
  \item Y. Yamaguchi, Y. Osamura, J.~D. Goddard, H.~F. Schaefer III,
        \emph{A New Dimension to Quantum Chemistry} (Oxford, 1994): the
        RHF and MCSCF second derivatives.
  \item N.~C. Handy, H.~F. Schaefer III, \emph{J.\ Chem.\ Phys.}\
        \textbf{81}, 5031 (1984): the $Z$-vector, and why it does not
        remove the $3N$ solves of a Hessian.
  \item P. Deglmann, F. Furche, R. Ahlrichs,
        \emph{Chem.\ Phys.\ Lett.}\ \textbf{362}, 511 (2002): an efficient
        RI SCF Hessian.
  \item R. Wang, Y. Wang, L. Lu, D. Hait, T.~J. Mart\'inez, ``Complete
        Active Space Self-Consistent Field with GPU-Accelerated Density
        Fitting'', ChemRxiv (2025), doi:10.26434/chemrxiv-2025-nd968, and
        its Supporting Information: generalized-density DF-J/K, SVD
        pseudo-coefficients, DF gradient contractions.
  \item J.~W. Snyder, E.~G. Hohenstein, N. Luehr, T.~J. Mart\'inez,
        \emph{J.\ Chem.\ Phys.}\ \textbf{143} (2015): AO-based SA-CASSCF
        gradients and the CP-MCSCF densities $D_{CV}$, $D_{\mathrm{Gen}}$
        (Ref.~52 of Wang et al.).
  \item T. Helgaker, P. J{\o}rgensen, \emph{Adv.\ Quantum Chem.}\
        \textbf{19}, 183 (1988): geometrical derivatives of variational
        wavefunctions, the connection, and the $2n+1$ rule.
  \item Q. Sun et al., \emph{WIREs Comput.\ Mol.\ Sci.}\ \textbf{8}, e1340
        (2018): PySCF, the reference throughout.
\end{itemize}

\end{document}
